\chapter{Esercizi}
\textbf{Es. 1}\\
Calcolare la lunghezza d'onda di un elettrone con energia cinetica $K_\alpha = \qty{6}{\mega\electronvolt}$.\\
\textbf{Soluzione:}\\
La lunghezza d'onda di un elettrone con energia cinetica $K_\alpha = \qty{6}{\mega\electronvolt}$ può essere calcolata utilizzando la relazione di de Broglie:
\begin{equation}
\lambda = \frac{h}{p}
\end{equation}
Otteniamo la quantità di moto $p$ utilizzando l'approssimazione relativistica $K_\alpha \simeq E_\alpha$ e dunque:
\begin{equation}
p = \sqrt{2m_e K_\alpha}
\end{equation}
dove $m_e$ è la massa dell'elettrone. Sostituendo $p$ nella formula della lunghezza d'onda:
\begin{equation}
\lambda = \frac{h}{\sqrt{2m_e K_\alpha}}
\end{equation}
Con i seguenti valori:
\begin{itemize}
 \item $h \simeq \qty{6.626e-34}{\joule\second}$
 \item $m_e \simeq \qty{9.109e-31}{\kilo\gram}$
 \item $K_\alpha$ espresso in joule (convertendo da MeV con $\qty{1}{\mega\electronvolt} = \qty{1.602e-13}{\joule}$)
\end{itemize}
Sostituendo i valori numerici, si ottiene $\lambda \simeq \qty{0.05}{\angstrom}$.\\
Utilizzando la formula semiempirica di massa e assumendo che le forze nucleari siano identiche per tutte le coppie di
nucleoni, possiamo scrivere l'energia di legame per nucleone come:\\
\textbf{Es. 2}\\
Fascio di particelle $\alpha$ con una corrente $i_\alpha = \qty{10}{\nano\ampere}$ colpisce un bersaglio di $^{24}\text{Mg}$ ($23.985 \,u$)con $\rho d = \qty{1}{\milli\gram\per\centi\meter^2}$. Creando una reazione isotropa:
\begin{equation}
^{24}_{12}\text{Mg} + ^4_2\alpha \rightarrow ^{27}_{13}\text{Al} + p
\end{equation}
Otteniamo un angolo solido di $\Delta \Omega = \qty{2e-4}{\steradian}$ con un numero di protoni al secondo $\dot{N}_p = \qty{20}{\second^{-1}}$. Calcolare la sezione d'urto differenziale totale.\\
\textbf{Soluzione:}\\
Sappiamo il numero di protoni che colpiscono il bersaglio per unità di tempo nel rivelatore posizionato ad un angolo solido $\Delta \Omega$, dunque otteniamo il numero totale di protoni prodotti per unità di tempo come:
\begin{equation}
  \dot{N}_{p,tot} = \dot{N}_{p}(\Delta \Omega) \dfrac{4 \pi}{\Delta \Omega} \simeq \qty{1.26e6}{\second^{-1}}
\end{equation}
Vogliamo trovare il flusso di particelle $\alpha$ che colpiscono il bersaglio:
\begin{equation}
 \phi_\alpha = \dfrac{i_\alpha}{Aq_\alpha}
\end{equation}
Il numero di centri di scattering è dato dai possibili centri di scattering del bersaglio di alluminio.
Possiamo anche scrivere $N_{\text{Mg}}$ in funzione della densità fornita:
\begin{equation}
 \dfrac{m_{\text{Mg}}N_{\text{Mg}}}{dA}d = \rho d
\end{equation}
Dunque abbiamo che:
\begin{equation}
  \sigma = \dfrac{\dot{N}_{p,tot}}{\phi_\alpha N_{\text{Mg}}} = \dfrac{\dot{N}_{p,tot}}{\dfrac{i_\alpha}{\cancel{A}q_\alpha} \dfrac{\rho d \cancel{A}}{m_{\text{Mg}}}} = \dfrac{\dot{N}_{p,tot} m_{\text{Mg}}q_\alpha}{i_\alpha \rho d A} \simeq \qty{160}{\milli\barn}
\end{equation}
\textbf{Es. 3}\\
Un elettrone con quantità di moto $p_-$ urta un positrone che si muove con quantità di moto $p_+$ uguale ed opposta. Nell'urto vengono prodotti una coppia di muoni con carica opposta:
\begin{equation}
 e^- + e^+ \rightarrow \mu^- + \mu^+
\end{equation}
Note le masse:
\begin{itemize}
 \item $m_e = \qty{0.511}{\mega\electronvolt\per c^}$
 \item $m_\mu = \qty{105.7}{\mega\electronvolt\per c^2}$
\end{itemize}
Determinare:
\begin{itemize}
  \item Le quantità di moto dei muoni in funzione di quelle degli elettroni
  \item Il valore minimo della quantità di moto degli elettroni affinché la reazione avvenga
\end{itemize}
\textbf{Soluzione:}\\
Utilizzando la conservazione della quantità di moto e dell'energia, possiamo scrivere:
\begin{equation}
 p_- + p_+ = p_{\mu^-} + p_{\mu^+} = 0
\end{equation}
\begin{equation}
 E_- + E_+ = E_{\mu^-} + E_{\mu^+}
\end{equation}
Dove:
\begin{equation}
 E_- = \sqrt{p_-^2 c^2 + m_e^2 c^4}
\end{equation}
\begin{equation}
 E_+ = \sqrt{p_+^2 c^2 + m_e^2 c^4}
\end{equation}
Dunque $E_- = E_+$ e $E_{\mu^-} = E_{\mu^+}$, quindi:
\begin{equation}
 \begin{split}
  2E_{e^-} &= 2E_{\mu^-} \\[8pt]
  E_{e^-} &= E_{\mu^-}
 \end{split}
\end{equation}
Sapendo che:
\begin{equation}
 E_{\mu^-} = \sqrt{p_{\mu^-}^2 c^2 + m_\mu^2 c^4}
\end{equation}
\begin{equation}
 E_{\mu^+} = \sqrt{p_{\mu^+}^2 c^2 + m_\mu^2 c^4}
\end{equation}
Sostituendo le quantità di moto opposte degli elettroni e risolvendo il sistema di equazioni, otteniamo:
\begin{equation}
 p_{\mu^\pm} = \sqrt{\dfrac{E_{e^-}^2}{c^2} - m_\mu^2 c^2} = \sqrt{p_{e^-}^2 + (m_e^2 - m_\mu^2) c^2}
\end{equation}
Otteniamo la quantità di moto minima iniziale ponendo $p_{\mu^\pm} = 0$:
\begin{equation}
 p_{e^-,min} = c \sqrt{m_\mu^2 - m_e^2} \simeq \qty{105.5}{\mega\electronvolt\per c}
\end{equation}
\section{Modello a gas di Fermi}
\textbf{Es. 1}\\
Calcolare l'impulso e l'energia di Fermi dei nucleoni nel nucleo $^{16}_8O$ assumendo una distribuzione a simmetria sferica. L'energia di legame totale del nucleo è $E_B = \qty{128}{\mega\electronvolt}$.\\
Calcolare la profondità della buca di potenziale che lega i nucleoni nel nucleo. (Si trascuri la differenza di massa tra protone e neutrone $M_P \simeq M_n = \qty{939}{\mega\electronvolt}$)\\
\textbf{Soluzione:}\\
Abbiamo che:
\begin{equation}
 p^n_F = p^p_F = p_F = \hbar \brackets{3 \pi^2 \dfrac{N}{V}}^{1/3}
\end{equation}
Dove $N = 8$ e $V = \dfrac{4}{3} \pi R^3$ con $R = r_0 A^{1/3}$ e $r_0 \simeq \qty{1.2}{\femto\meter}$. Dunque:
\begin{equation}
 p_F = \hbar \brackets{3 \pi^2 \dfrac{N}{\dfrac{4}{3} \pi r_0^3 A}}^{1/3} \simeq \qty{250}{\mega\electronvolt\per c}
\end{equation}
Ora troviamo l'energia di Fermi:
\begin{equation}
 \varepsilon_F = \sqrt{p_F^2 c^2 + M_N^2 c^4} - M_N c^2 \simeq \qty{33}{\mega\electronvolt}
\end{equation}
Per calcolare la profondità della buca di potenziale sappiamo che:
\begin{equation}
  u = \varepsilon_F + |E_B|/A \simeq \qty{41}{\mega\electronvolt}
\end{equation}
\section{Deutoni}
\textbf{Es. 1}\\
Un fascio di deutoni (D), caratterizzato da una corrente $i_D = \qty{2}{\micro\ampere}$, colpisce un bersaglio di trizio solido con spessore superficiale $\rho d = \qty{0.2}{\milli\gram\per\centi\meter^2}$. Considerando la reazione:
\begin{equation}
 ^3_1\text{H} + ^2_1\text{D} \rightarrow ^4_2\text{He} + n
\end{equation}
Assumendo di avere un ricevoitore con superficie pari a $S=\qty{20}{\centi\meter^2}$ posto ad una distanza $r=\qty{3}{\meter}$ dal bersaglio, calcolare il numero di neutroni al secondo rilevati, sapendo che la sezione d'urto differenziale della reazione è $\dv{\sigma}{\Omega} = \qty{13}{\milli\barn\per\steradian}$ ad un angolo $\theta = 30^\circ$.\\
\textbf{Soluzione:}\\
Calcoliamo il flusso di deutoni che colpisce il bersaglio:
\begin{equation}
 \phi_D = \dfrac{i_D}{q_D S} = \dfrac{i_D}{2e S}
\end{equation}
Il numero di nuclei di trizio nel bersaglio è dato da:
\begin{equation}
 N_{^3\text{H}} = \dfrac{\rho d S N_A}{M_{^3\text{H}}}
\end{equation}
Dove $N_A$ è il numero di Avogadro e $M_{^3\text{H}} \simeq \qty{3}{\gram\per\mole}$ è la massa molare del trizio.\\
Poichè l'angolo solido $\Delta \Omega$ è dato da:
\begin{equation}
 \Delta \Omega = \dfrac{S}{r^2}
\end{equation}
e risulta $\Delta \Omega \rightarrow 0 \implies \dv{\sigma}{\Omega} \Delta \Omega \simeq \sigma$, possiamo scrivere:
\begin{equation}
  \dot{N}_n = \phi_D N_{^3\text{H}} \dv{\sigma}{\Omega} \Delta \Omega
\end{equation}
Sostituendo i valori numerici, otteniamo:
\begin{equation}
 \dot{N}_n \simeq \qty{1.3e4}{\second^{-1}}
\end{equation}

\section{Prove d'esame}
\subsection{Prova 1}
\textbf{Es. 1}\\
Supponiamo di osservare sperimentalmente il decadimento di una particella di massa $M$ in due fotoni. L'angolo di apertura tra i due fotoni emessi è pari a $\vartheta = 32^\circ$ e le energie dei due fotoni
sono misurate essere $K_1 = \qty{200}{\mega\electronvolt}$ e $K_2 = \qty{300}{\mega\electronvolt}$, rispettivamente. Determinare la massa $M$ della particella originaria e la sua quantità di moto nel sistema del laboratorio prima del decadimento.\\
\textbf{Soluzione:}\\
Utilizzando la conservazione dell'energia possiamo scrivere:
\begin{equation}
    E = K_1 + K_2 = \qty{500}{\mega\electronvolt}
\end{equation}
L'energia iniziale della particella decaduta è data dalla relazione di dispersione relativistica:
\begin{equation}
    E^2 = p^2 c^2 + M^2 c^4
\end{equation}
Scegliamo l'asse $x$ lungo la direzione della particella decaduta, dunque:
\begin{equation}
    p_x = p_{1} \cos(\theta_1) + p_{2} \cos(\theta_2) = P
\end{equation}
Dalla conservazione della quantità di moto sappiamo:
\begin{equation}
    p_y = p_{1} \sin(\theta_1) - p_{2} \sin(\theta_2) = 0
\end{equation}
La quantità di moto dei fotoni è data da:
\begin{equation}
    \begin{split}
        p_{1} = \dfrac{K_{1}}{c} \\[8pt]
        p_{2} = \dfrac{K_{2}}{c}
    \end{split}
\end{equation}
Dunque otteniamo:
\begin{equation}
    \dfrac{\sin(\theta_1)}{\sin(\theta_2)} = \dfrac{p_{2}}{p_{1}} = \dfrac{K_{2}}{K_{1}} = 1.5
\end{equation}
Sapendo che $\vartheta = \theta_1 + \theta_2 = 32^\circ$, tramite le formule di addizione e sottrazione esprimiamo tutto in funzione di $\theta_2$:
\begin{equation}
    \sin(\theta_1) = \sin(32^\circ - \theta_2) = \sin(32^\circ) \cos(\theta_2) - \cos(32^\circ) \sin(\theta_2)
\end{equation}
Da cui:
\begin{equation}
    \tan(\theta_2) = \dfrac{\sin(32^\circ)}{1.5 + \cos(32^\circ)} \simeq 0.183
\end{equation}
Dunque $\theta_2 \simeq 12.7^\circ$ e $\theta_1 \simeq 19.3^\circ$. Ora possiamo calcolare la quantità di moto della particella decaduta:
\begin{equation}
    P = \dfrac{K_{1}}{c} \cos(19.3^\circ) + \dfrac{K_{2}}{c} \cos(12.7^\circ) \simeq \qty{481}{\mega\electronvolt\per c}
\end{equation}
Ora possiamo calcolare la massa della particella decaduta:
\begin{equation}
    M = \dfrac{1}{c^2} \sqrt{E^2 - P^2 c^2} \simeq \qty{137}{\mega\electronvolt\per c^2}
\end{equation}
La massa ricavata è compatibile con quella del mesone $\pi^0$.\\
\textbf{Es. 2}\\
Calcolare la dose e la dose equivalente assorbita in 4 ore da una persona di $M_p = \qty{100}{\kilo\gram}$ che ha ingerito $\qty{0.1}{\gram}$ di Americio $\atom{Am}{241}{95}$ (peso atomico 243 u), nuclide che decade emettendo una particella $\alpha$ di energia cinetica $K_\alpha = \qty{5.6}{\mega\electronvolt}$ con tempo di dimezzamento pari a 432.2 anni (assumere $FQ=10$).\\
\textbf{Soluzione:}\\
Il numero di atomi di Americio ingeriti è dato da:
\begin{equation}
    N_{\text{Am}} = \dfrac{m_{\text{Am}} N_A}{M_{\text{Am}}} \simeq 2.48 \cdot 10^{20}\,\text{atomi}
\end{equation}
Grazie al tempo di dimezzamento possiamo calcolare la costante di decadimento:
\begin{equation}
    \lambda = \dfrac{\ln(2)}{T_{1/2}} \simeq \qty{5.08e-11}{\per\second}
\end{equation}
Convertendo le 4 ore in secondi, otteniamo il numero di decadimenti avvenuti nel tempo considerato:
\begin{equation}
    N_d(\Delta t) = N_{\text{Am}} (1 - \efunction^{-\lambda \Delta t}) \simeq 1.81 \cdot 10^{14}\,\text{decadimenti}
\end{equation}
L'energia totale assorbita dalla persona è data da:
\begin{equation}
    E_{\text{abs}} = N_d(\Delta t) K_\alpha \simeq \qty{10.14e14}{\mega\electronvolt} = \qty{162}{\joule}
\end{equation}
La dose assorbita è data da:
\begin{equation}
    D = \dfrac{E_{\text{abs}}}{M_p} \simeq \qty{1.62}{\gray}
\end{equation}
La dose equivalente è data da:
\begin{equation}
    H = FQ \cdot D \simeq \qty{16.2}{\sievert}
\end{equation}
\subsection{Prova 2}
\textbf{Es. 1}\\
Il trizio $\atom{H}{3}{1}$ ($M_t$ = 3 u) decade $\beta$ con vita media $\tau = 17.77$ anni. Sapendo che un campione arricchito di idrogeno gassoso, contenente $\qty{0.1}{\gram}$ di trizio, produce 21 calorie/ora, stimare l'energia media delle particelle $\beta$ emesse.\\
\textbf{Soluzione:}\\
Il numero di atomi di trizio nel campione è dato da:
\begin{equation}
    N_{\text{H-3}} = \dfrac{m_{\text{H-3}} N_A}{M_{\text{H-3}}} \simeq 2.01 \cdot 10^{22}\,\text{atomi}
\end{equation}
In un'ora, il numero di atomi di trizio che decadono è dato da:
\begin{equation}
    N_d(\Delta t) = N_{\text{H-3}} (1 - \efunction^{-\Delta t/\tau}) \simeq 1.29 \cdot 10^{17}\,\text{decadimenti}
\end{equation}
Per ogni decadimento viene emessa una particella $\beta$, dunque l'energia totale emessa in un'ora è:
\begin{equation}
    E_{\text{tot}} = 21\,\text{calorie} = \qty{87.96}{\joule}
\end{equation}
Dunque l'energia media delle particelle $\beta$ emesse è:
\begin{equation}
    \avg{K}_\beta = \dfrac{E_{\text{tot}}}{N_d} \simeq \qty{6.82e-16}{\joule} = \qty{4.26}{\kilo\electronvolt}
\end{equation}
Alternativamente sappiamo che il numero di decadimente per unità di tempo è dato da:
\begin{equation}
    \dot{N}_d = \lambda N_{\text{H-3}}
\end{equation}
Dunque l'energia media delle particelle $\beta$ emesse può essere calcolata come:
\begin{equation}
    \avg{K}_\beta = \dfrac{P}{\dot{N}_d} = \dfrac{P}{\lambda N_{\text{H-3}}} \simeq \qty{4.26}{\kilo\electronvolt}
\end{equation}
